\begin{abstract}
Urban last-mile logistics increasingly relies on mixed commercial fleets composed of battery electric vehicles (BEVs), internal combustion engine (ICE) vans, and hybrid powertrains. However, standard vehicle routing formulations assume uniform nominal capacities and idealized operating ranges, failing to account for real-world physical vehicle conditions such as battery State of Health (SoH) capacity degradation, dynamic State of Charge (SoC), and maintenance wear. Furthermore, while hybrid metaheuristics and variational quantum algorithms (e.g., Quantum Approximate Optimization Algorithm, QAOA) have been proposed, existing hybrid routing systems are typically configured as unidirectional linear pipelines or isolated candidate comparisons without genuine algorithmic cooperation, bidirectional candidate exchange, or online learning.

To address these fundamental challenges, this paper presents \textbf{ROUTE/IQ-NextGen}, an adaptive, vehicle-condition-aware cooperative optimization framework. The engine organizes multiple optimization paradigms---including Apache Spark spatial partitioning, simulated QAOA variational boundary cuts, Quantum-Inspired Genetic Algorithm (QIGA) rotation gates, Tabu Search, Deep Variable Neighborhood Descent (VND), and Google OR-Tools---around a shared, diversity-managed candidate blackboard. The framework introduces four primary research contributions: (1) an explicit \textbf{Heterogeneous Fleet Vehicle Condition (HFVRP-VC)} formulation that bounds routes by usable battery energy ($E = B_{\text{nom}} \cdot \text{SoH} \cdot \text{SoC}$) and enforces multi-objective trade-offs across distance, operating cost, energy, and carbon emissions; (2) a \textbf{Topological Edge Jaccard Diversity Metric} that prevents candidate pool stagnation; (3) an \textbf{Adaptive Contextual Multi-Armed Bandit (MAB)} controller with Upper Confidence Bound (UCB-1) learning and a multi-objective feedback reward penalizing execution runtime while rewarding feasibility and diversity; and (4) a proven \textbf{Conditional Non-Degradation Guardrail (Proposition 1)} guaranteeing by construction that final fleet distance never exceeds the classical OR-Tools baseline.

Empirically validated on real-world New York City delivery networks across scales $N \in [12, 30]$ with heterogeneous vehicle fleets, ROUTE/IQ-NextGen achieves a mean distance reduction of \textbf{$-1.24\%$} over Google OR-Tools alone (reaching \textbf{$-3.16\%$} on $N=30$ nodes), while reducing mean execution latency by \textbf{$24.0\%$} ($1219.2\text{ ms}$ vs $1604.0\text{ ms}$) compared to legacy cooperative architectures through adaptive operator selection. Ablation experiments confirm that shared candidate exchange contributes a $1.07\%$ distance improvement over isolated waterfall execution, while QIGA and QAOA operators generate high topological diversity ($\mathcal{D} = 0.41\text{--}0.61$). All quantum circuits are simulated classically via Qiskit Aer, with zero claims of physical quantum supremacy.
\end{abstract}

\begin{IEEEkeywords}
Cooperative optimization, heterogeneous fleet vehicle routing, vehicle condition awareness, battery degradation, contextual multi-armed bandit, QAOA, QIGA, Apache Spark.
\end{IEEEkeywords}
